\chapter{Introduction}
\label{ch:intro}

\section{Research background}
\label{sec:intro-background}

Customer segmentation converts behavioural records into groups that can be inspected, compared and queried. A clustering algorithm never observes a customer directly. It receives a numeric vector assembled from selected events, time windows and aggregation rules — this vector, and the set of rules used to build it, is what this dissertation calls a customer \emph{representation}. A compact purchase summary emphasises transaction timing and repetition. A richer representation also records browsing volume, conversion ratios, temporal engagement and category breadth. Because a clustering algorithm only ever sees the representation, not the customer, changing the representation changes the distances the algorithm measures between the same people, and can make different group structures visible or invisible.

This dissertation studies that effect for one specific, bounded population: visitors to an e-commerce platform who completed at least one purchase, drawn from the RetailRocket event log described in Chapter~\ref{ch:methodology}. The general problem — that a clustering algorithm's output depends on how customers are encoded, not only on which customers are studied — is not specific to this dataset. The empirical evidence that follows, however, is: it demonstrates the effect for this purchasing cohort under one feature contract, and does not claim that the same magnitude of effect, or even the same direction of effect, would be observed for non-purchasing visitors, a different platform, or a different set of added variables.

Representation is often evaluated only indirectly, through the score of whatever clustering the analyst finally settles on. That approach cannot separate the effect of the information choice from the effect of algorithm behaviour, because the cohort, the parameter grid and the validation rule are usually all changing at the same time. It can also reward a clustering with a high silhouette score (an internal measure of how compact and well separated clusters are, defined precisely in \S\ref{sec:candidates}) that isolates a handful of unusual observations while leaving almost every other customer in a single undivided group — a partition that scores well but is operationally useless. \emph{Clusterability} research — the study of whether, and by which criterion, a dataset can be said to contain genuine recoverable group structure at all — shows that this question depends on the definition adopted and on the data's own conditions rather than on one universal test \citep{ackerman2009clusterability,adolfsson2019clusterability}. Representation must therefore be examined as part of the empirical model being fitted to customer event logs, not treated as a preliminary, inconsequential step that precedes the "real" modelling choices.

Explanation introduces a second, independent source of ambiguity. CLASSIX, one of the two clustering methods used in this dissertation, produces \emph{provenance}: a traceable, step-by-step record of which points were grouped together, in what order, and by which rule, so that a customer's final cluster label can be traced back through the intermediate groups that produced it \citep{chen2024classix}. A decision tree fitted afterwards to approximate K Means's cluster labels — the approach taken by ExKMC in this dissertation — instead provides a \emph{post hoc} approximation, whose accuracy against the original labels (its \emph{fidelity}) and whose size (its rule count) can both be measured \citep{moshkovitz2020imm}. These two kinds of explanation answer different questions about different objects. A short, high-fidelity tree can closely reproduce a reference partition without revealing anything about how that partition was originally computed. A detailed intrinsic trace can account for every single assignment while still being too large to summarise as a handful of business rules. Treating these two kinds of evidence as interchangeable, or combining them into one explainability score, would hide exactly the effect this dissertation is designed to detect: that representation changes both the cluster structure itself and the burden of explaining it.

\section{Research gap}
\label{sec:intro-gap}

Two gaps in the existing literature motivate this study, and both are developed further in Chapter~\ref{ch:litreview}. First, representation is usually held fixed, or varied only incidentally, while attention is given to algorithm choice or to the final clustering score; the joint, controlled effect of representation on both structure and stability, for the same customers under matched grids and constraints, is rarely isolated directly. Second, explainable clustering methods are frequently discussed as a single category, when they in fact occupy at least three distinct design positions: \emph{intrinsic} methods (such as CLASSIX) that expose the computation producing their own partition; \emph{explainability-first} methods that build a simple description into the clustering objective itself; and \emph{post hoc} methods (such as ExKMC) that approximate an already-fitted reference partition afterwards. These three positions are not interchangeable, and conflating them risks attributing a property demonstrated for one (for example, ExKMC's measured fidelity) to a different kind of evidence (for example, CLASSIX's own provenance).

\section{Research questions}
\label{sec:intro-rq}

This dissertation answers five research questions, each of which corresponds directly to one part of the experimental programme carried out in Chapters 3 and 4:

\begin{enumerate}[label=RQ\arabic*.]
\item How much do independently selected K Means and CLASSIX partitions of the same purchasers agree with each other when the customers are described by a compact transaction representation compared with a richer behavioural representation?
\item Does high within-representation resampling stability (repeatability of a fitted configuration when refit on repeated subsamples of the same representation) coexist with low cross-representation agreement, and what does that combination imply about representation as a modelling choice rather than a preprocessing step?
\item How does representation change the two available explanation channels — CLASSIX's own intrinsic provenance and ExKMC's post hoc, held-out fidelity — and the computational or complexity profile of each?
\item Which selected clusters, if any, are operationally addressable (large enough, sufficiently data-complete, and distinctive enough on named features to be targeted with a simple query), and how does this depend on representation and clustering method?
\item Do the observed representation effects persist under sensitivity checks — feature deletion, PCA, small numerical perturbations, alternative session-window definitions, and a permutation control — or are they artefacts of one specific parameter choice?
\end{enumerate}

\section{Research purpose, significance and contribution}
\label{sec:intro-purpose}

The purpose of this dissertation is to answer these five questions for one matched pair of representations of the same 11,719 RetailRocket purchasers. The compact representation contains three variables describing purchase recency and frequency. The richer representation retains those three variables unchanged and adds nine further variables describing activity volume, conversion, temporal engagement and category behaviour. Both representations use the same event snapshot, the same customer order, the same transformations, the same algorithm grids, the same deployment constraints and the same resampling seeds. This design compares the effect of adding behavioural information within one recorded population; it does not compare different cohorts, different snapshots, or different algorithms' default settings against each other.

Holding the cohort and the evaluation rules fixed in this way is what makes the comparison informative: it isolates the effect of representation from changes in population or in model selection that would otherwise be entangled with it. The analysis also keeps CLASSIX's intrinsic provenance and ExKMC's post hoc approximation strictly separate, because — as argued above — they explain different objects and are not comparable on one shared scale. A supporting mathematical section (\S\ref{sec:methods}) restates why CLASSIX's projection-based aggregation search can safely skip some full-distance comparisons (its \emph{projection filtering} step), and derives a local, sufficient condition under which that aggregation search provably reaches the same result before and after a small perturbation of the representation.

This dissertation makes five contributions, corresponding to RQ1–RQ5 above:

\begin{enumerate}
\item a paired comparison of compact and rich representations of the same customers, rather than a comparison across unrelated cohorts or studies;
\item a joint clustering-structure and resampling-stability comparison for two different algorithm families (K Means and CLASSIX) under matched grids and deployment constraints;
\item an explicit separation of CLASSIX's intrinsic provenance from ExKMC's post hoc fidelity as two distinct, non-interchangeable kinds of explanation evidence;
\item a multi-part robustness and sensitivity analysis (feature deletion, PCA, numerical perturbation, alternative session windows, and a permutation control) that tests which parts of the representation effect survive specific, named alternative explanations;
\item a bounded mathematical analysis of CLASSIX's projection-filtering step and of the conditions under which its distance-merging aggregation result is provably stable to small perturbations.
\end{enumerate}

None of these contributions supports, and none is intended to support, a claim that either representation or either clustering method is universally superior; each is scoped to what the corresponding evidence in Chapter~\ref{ch:results} actually establishes.

\section{Research content and headline findings}
\label{sec:intro-content}

The empirical programme begins with nine externally labelled benchmark datasets, used to check that the four clustering implementations used later — K Means, CLASSIX, and, for this benchmark only, two comparison methods, DBSCAN (a density-based clustering algorithm) and Ward's method (an agglomerative hierarchical clustering algorithm) — and the label-free selection procedure behave sensibly across a range of known cluster shapes before either is applied to the RetailRocket data. DBSCAN and Ward's method appear only in this supporting benchmark; they are not used anywhere in the main RetailRocket analysis. The main analysis then applies K Means and CLASSIX to both RetailRocket representations and compares cluster counts, internal indices, customer assignments and fixed-parameter overlap-resampling stability (the repeatability of a fitted configuration across repeated 80 per cent subsamples). CLASSIX provenance and held-out ExKMC trees provide the two explanation channels described above. Feature deletion, PCA, numerical perturbation, alternative session windows and a permutation control examine sensitivity. A supporting mathematical analysis (\S\ref{sec:projection-filtering}--\S\ref{sec:stability-certificate}) states the scope of CLASSIX's projection filtering, its distance complexity, the conditions for deterministic execution, and a local sufficient stability condition.

The headline empirical result concerns cluster agreement. The two independently selected K Means partitions — one fitted on the compact representation, one on the rich representation, of the \emph{same} 11,719 customers — agree with each other at an adjusted Rand index (ARI) of 0.014; the two CLASSIX partitions agree at 0.068. The adjusted Rand index measures agreement between two partitions of the same objects: it equals 1 for identical partitions and is corrected so that pairing labels at random scores approximately 0. Critically, there is no external ground-truth segmentation of these customers to compare against — the "two partitions" being compared throughout this dissertation are always the compact-space and rich-space clusterings of the same people, never a clustering against a known correct answer. An ARI of 0.014 or 0.068 therefore does not mean either clustering is 1.4\% or 6.8\% "accurate"; it means the two representations of the same customers produce partitions that are little more similar to each other than two unrelated random groupings would be, despite both being computed from the same underlying event history. Compact representations give higher silhouette values, while the rich K Means partition is more balanced across its four clusters and the rich CLASSIX partition requires more internal aggregation groups (a measure of computational provenance, not necessarily of human interpretability). At the common four-leaf ExKMC budget, held-out fidelity is 0.9448 for the compact K Means approximation and 0.9650 for the rich one. Of the resulting clusters, only one — a rich-representation K Means cluster — passes the documented addressability audit. Together, these findings show that changing the represented information changes both the observed grouping and the burden of explaining it, for this cohort and this feature contract; they do not establish that one representation is universally preferable, nor that any resulting cluster has realised business value.
